\documentclass[11pt]{article}

\usepackage{extras} % Se extras.sty
\usepackage{tikz}
\usepackage{parskip}
\usepackage{standalone}
%\usepackage{titlesec}
\usetikzlibrary{positioning}
\usetikzlibrary{shapes}
\usetikzlibrary{calc}
\usepackage{enumitem}
\def\arraystretch{1.5}
\graphicspath{ {images/} }
\setlength{\LTpost}{0pt}
\pagenumbering{roman}
\setlength{\LTpost}{0pt}%
\usepackage{graphicx}
\usetikzlibrary{patterns}
\usetikzlibrary{shapes.arrows}
\usepackage{amsfonts}
\usetikzlibrary{positioning}
\usetikzlibrary{calc}
\usetikzlibrary{decorations.pathmorphing}
\usepackage{amssymb}
\usetikzlibrary{shapes,arrows}

%\usepackage[mode=buildnew]{standalone}

%\titlespacing*{\subsection}{0pt}{1.1\baselineskip}{\baselineskip}


\begin{document}
\begin{titlepage}
\begin{center}

{\Large\bfseries TSEA56 - Kandidatprojekt i elektronik \\ LIPS Systemskiss}

\vspace{5em}

Version 0.1

\vspace{5em}
%
Grupp 4 \\
\begin{tabular}{rl}
Hynén Ulfsjöö, Olle&\verb+ollul666+
\\
Wasteson, Emil&\verb+emiwa068+
\\
Tronje, Elena&\verb+eletr654+
\\
Gustafsson, Lovisa&\verb+lovgu777+
\\
Inge, Zimon&\verb+zimin415+
\\
Strömberg, Isak&\verb+isast763+
\\
\end{tabular}

\vspace{5em}
\today

\vspace{16em}
Status
\begin{longtable}{|l|l|l|} \hline

Granskad & - & - \\ \hline
Godkänd & - & - \\ \hline
 
\end{longtable}

\end{center}
\end{titlepage}

\pagebreak
\begin{center}

\section*{PROJEKTIDENTITET}
2016/VT, Undsättningsrobot Gr. 4

Linköpings JÄTTEtekniska högskola, ISY
\vspace{5em}
\begin{center}

\begin{tabular}{|l|l|l|l|} \hline
\textbf{Namn} & \textbf{Ansvar} & \textbf{Telefon} & \textbf{E-post}  \\ \hline 
Isak Strömberg (IS) & Projektledare & 073-980 38 50 & isast763@student.liu.se \\ \hline
Olle Hynén Ulfsjöö (OHU)& Dokumentansvarig & 070-072 91 84 & ollul666@student.liu.se \\ \hline
Emil Wasteson (EW) & Hårdvaruansvarig & 076-836 61 66 & emiwa068@student.liu.se \\ \hline
Elena Tronje (ET) & Mjukvaruansvarig & 072-276 92 93 & eletr654@student.liu.se \\ \hline
Zimon Inge (ZI)& Testansvarig & 070-171 35 18 & zimin415@student.liu.se \\ \hline
Lovisa Gustafsson (LG) & Leveransansvarig & 070-210 32 53 & lovgu777@student.liu.se \\ \hline
\end{tabular}

\end{center}

E-postlista för hela gruppen: isast763@student.liu.se

\vspace{5em}
Kund: ISY, Linköpings universitet \\
tel: 013-28 10 00, fax: 013-13 92 82 \\
Kontaktperson hos kund: Mattias Krysander \\
tel: 013-28 21 98, e-post: matkr@isy.liu.se \\

\vspace{5em}
Kursansvarig:  Tomas Svensson\\
tel: 013-28 13 68, e-post: tomass@isy.liu.se \\
Handledare: Peter Johansson \\
tel: 013-28 13 45, epost: peter.a.johansson@liu.se
\end{center}
\pagebreak

\tableofcontents

\pagebreak

\section*{Dokumenthistorik}
\begin{table}[h]
\begin{tabular}{|l|l|l|l|l|} \hline

\textbf{Version} & \textbf{Datum} & \textbf{Utförda förändringar} & \textbf{Utförda av} & \textbf{Granskad} \\ \hline
0.1 & FYLL I &  Första utkastet & Grupp 4 & - \\ \hline
\end{tabular}
\end{table}

\pagebreak
\pagenumbering{arabic}

\begin{flushleft}

\section{Inledning}
Denna systemskiss utgör en del av förarbetet i kandidatprojektet TSEA56. På uppdrag av beställaren ska systemskissen ge en detaljerad beskrivning ar produkten och dess delsystem. Dokumentet inleds med en grov beskrivning av produkten, därefter följer en mer detaljerad beskrivning på modulnivå.
\subsection{Syfte och mål}
Syftet med projektet är att bygga en undsättningsrobot på prototypnivå. Roboten ska autonomnt kunna utforska en \textit{simulerad} grotta och finna en nödställd. Parallellt med sökande ska roboten kunna kommunicera med en datormodul där en karta av grottan successivt ritas upp. När den nödställde är funnen ska en optimerande algoritm beräkna den kortaste vägen dit. Denna väg ska sedan användas när roboten ska förse ned nödställde med någon förnödenhet. Dessa krav finns specificerade i kravspecifikationen.\cite{krav}

\pagebreak
\section{Översikt av systemet}
Roboten i sin miljö finns illustrerad i figur \ref{system}. Kommunikationen med datormodulen sker åt båda hållen och via Bluetooth\textsuperscript{\circledR}. Roboten ska dock även klara sitt uppdrag utan kommunikation med datormodulen. Det vill säga att kartläggning, styrning och optimering av kortaste väg sker lokalt på roboten. Banan är uppbyggd enligt banspecifikationen \cite{banspec} och uppdraget utförs enligt tävlingsreglerna \cite{tavling}.
\begin{figure}[htbp]
\centering
\noindent\resizebox{.8\linewidth}{!}{
	
\documentclass[border=10px]{standalone}
\usepackage{tikz}
\usetikzlibrary{patterns}
\usetikzlibrary{shapes.arrows}
\usepackage{amssymb}
\usetikzlibrary{calc}
\usepackage{verbatim}
\usetikzlibrary{decorations.pathmorphing}

\begin{document}
	
\begin{tikzpicture}[scale=1,rotate=0]
		
	%Base
	\draw[thick, draw=black, fill=gray!10] (0,0) rectangle (6,10);

	%Wheels
	\draw[thick, pattern=north west lines, pattern color=black] (-.5,1) 		rectangle (0,2.5);
	\draw[thick, pattern=north west lines, pattern color=black] (-.5,7.5) 	rectangle (0,9);
	\draw[thick, pattern=north west lines, pattern color=black] (6,1) 		rectangle (6.5,2.5);
	\draw[thick, pattern=north west lines, pattern color=black] (6,7.5) 		rectangle (6.5,9);
	
	%Sensors
	\draw[thick, draw=black, fill=white]		(-.25,.25) 		rectangle 	(.5,.75);
	\draw[thick, draw=black, fill=white] 	(-.25,9.25) 		rectangle 	(.5,9.75);
	\draw[thick, draw=black, fill=white] 	(5.5,.25) 		rectangle 	(6.25,.75);
	\draw[thick, draw=black, fill=white] 	(5.5,9.25) 		rectangle 	(6.25,9.75);
	\draw[thick, draw=black, fill=white] 	(1.5,10.25) 		rectangle 	(4.5,9.5);
	\draw[thick, draw=black]				 	(2.5,9.5)			--		  	(2.5,10.25);
	\draw[thick, draw=black]					(3.5,9.5) --++ (0,.75);
	
	%Gripklo
	\draw[thick, draw=black, <-] (3,10.25) --++ (0,.72) node[above] {Gripklo};
	
	%Lidar lite
	\draw[thick, draw=black, <-] 			(1.5,10.25)  		--+ 		  	(135:1) node[above, align=center] {\verb+LIDAR-Lite v2+ \\ med servo};
	%Identifierare av nödställd
	\draw[thick, draw=black, <-]				(4.5,10.25)		--+			(45:1)  node[above] {\verb+IRM-8601-S+};
	
	%IR sensorer
	\draw[thick, draw=black, <-]				(-.25, 9.75) 	--+			(135:1) node[above] {\verb+GP2D120+};
	\draw[thick, draw=black, <-]				(-.25, .25) 		--+			(225:1) node[below] {\verb+GP2D120+};
	\draw[thick, draw=black, <-]				(6.25, 9.75) 	--+			(45:1) node[above] {\verb+GP2D120+};
	\draw[thick, draw=black, <-]				(6.25, .25) 		--+			(315:1) node[below] {\verb+GP2D120+};
	
	\draw[thick, draw=black, fill=white]					(2.5,4.5)		rectangle	(3.5,5.5);
	
	%Gyro
	\draw[thick, draw=black, <-] (3.5,5) --++ (.25,0) node[rotate=270, above] {\verb+MLX90609+};
	
	%Arrows and text
	%\draw[thick, ->]  (3,11) node[left, align=center] {\verb+LIDAR-Lite v2+ \\ + detektor av nödställd} -- (3,10.25);
	%\draw[thick, <->] (0.5,0.5)  --  (5.5,0.5) node[left=-14pt,midway, fill=gray!10] {\verb+GP2D120+};
	%\draw[thick, <->] (0.5,9.5) -- (3,9) node[right=-14pt,fill=gray!10] {\verb+GP2D120+} -- (5.5,9.5);
	%\draw[thick, ->] (4.5,5) node[above] {\verb+MPU-6500+} -- (3.5,5);
	
	%Sensormodul
	\node (sensormodul) at (3,6.75) [thick, draw=black, minimum width=3cm, minimum height=1.5cm, align=center, fill=white] {Sensormodul \\ \verb+ATmega1284p+};
	
	%Huvudmodul
	\node (huvudmodul) at (3,3.25) [thick, draw=black, minimum width=3cm, minimum height=1.5cm, align=center, fill=white] {Huvudmodul \\ \verb+ATmega1284p+};
	
	%Styrmodul
	\node (styrmodul) at (3,1.25) [thick, draw=black, minimum width=3cm, minimum height=1.5cm, align=center, fill=white] {Styrmodul \\ \verb+ATmega1284p+};
	
	%I2C-buss
	\draw[thick, draw=teal] ($(sensormodul.north east) + (.25,0)$) -- ($(styrmodul.south east) + (.25,0)$);
	\draw[thick, draw=teal] ($(sensormodul.north east) + (.5,0)$) -- ($(styrmodul.south east) + (.5,0)$);
	\node (i2c) at ($(sensormodul.north east) + (0.375,.25)$) [text=teal] {\small I\textsuperscript{2}C};
	
	\draw[thick, draw=teal, fill=teal] ($(sensormodul.south east) + (0,.25)$) --+ (.25,0) circle [radius=.05cm];
	\draw[thick, draw=teal, fill=teal] ($(sensormodul.south east) + (0,.5)$) --+ (.5,0) circle [radius=.05cm];
	
	\draw[thick, draw=teal, fill=teal] ($(huvudmodul.south east) + (0,.25)$) --+ (.25,0) circle [radius=.05cm];
	\draw[thick, draw=teal, fill=teal] ($(huvudmodul.south east) + (0,.5)$) --+ (.5,0) circle [radius=.05cm];
	
	\draw[thick, draw=teal, fill=teal] ($(styrmodul.south east) + (0,.25)$) --+ (.25,0) circle [radius=.05cm];
	\draw[thick, draw=teal, fill=teal] ($(styrmodul.south east) + (0,.5)$) --+ (.5,0) circle [radius=.05cm];
	
	%Koppling till sensormodul
	\draw[thick, draw=violet] (.5,9.5) --++ (.25,0) --++ (0,-2.5) --++ (0.75,0) ;
	\draw[thick, draw=violet] (.5,.5) --++ (.25,0) --++ (0,6.25) --++ (0.75,0);
	
	\draw[thick, draw=violet] (5.5,9.5) --++ (-.25,0) --++ (0,-2.5) --++ (-0.75,0);
	\draw[thick, draw=violet] (5.5,.5) --++ (-.25,0) --++ (0,6.25) --++ (-0.75,0);
	
	\draw[thick, draw=violet] (2,9.5) --++ (0,-2);
	\draw[thick, draw=violet] (4,9.5) --++ (0,-2);
	
	\draw[thick, draw=violet] (sensormodul.south) --++ (0,-.5);
	
	%Kopplingar till styrmodul
	\draw[thick, draw=olive] (styrmodul.east) --++ (1.5,0);
	\draw[thick, draw=olive] (styrmodul.west) --++ (-1.5,0);
	\draw[thick, draw=olive] ($(styrmodul.east) + (0,.25)$) --++ (1,0) --++ (0,6.25) --++ (.5,0);
	\draw[thick, draw=olive] ($(styrmodul.west) + (0,.25)$) --++ (-1,0) --++ (0,6.25) --++ (-.5,0);
	\draw[thick, draw=olive] (3,9.5) --++ (0,-1.75) --++ (-1.75,0) --++ (0,-6) --++ (.25,0);
	
	%Bluetooth
	\draw[thick, draw=black] (huvudmodul.east) --++ (2,0) node[right,minimum width=.75cm, minimum height=1cm, draw=black] (bt) {\includegraphics{bluetooth}};
	
	\node (bt2) at ($(bt) + (2,0)$) [thick, draw=black, minimum width=.75cm, minimum height=1cm] {\includegraphics{bluetooth}};
	
	\draw[thick, ->,line join=round,decorate, decoration={
    												snake,
    												segment length=5,
    												amplitude=1,
    												post=lineto,
    												post length=1pt}] 
    		($(bt.east) + (5pt,5pt)$) -- ($(bt2.west) + (-5pt,5pt)$);
    		
    \draw[thick, ->,line join=round,decorate, decoration={
    												snake,
    												segment length=5,
    												amplitude=1,
    												post=lineto,
    												post length=1pt}] 
    		 ($(bt2.west) + (-5pt,-5pt)$) -- ($(bt.east) + (5pt,-5pt)$);
    		 
    \draw[thick, draw=black] (bt2.east) --++ (.5,0) node[right, minimum width=3cm, minimum height=1.5cm, draw=black, fill=white] {Datormodul};
    
    \draw[thick, draw=black] ($(sensormodul.south west) + (.25,0)$) --++ (0,-2) node[midway, above, sloped] {\small avbrott};
    
    \draw[thick, draw=black] ($(styrmodul.north west) + (.25,0)$) -- ($(huvudmodul.south west) + (.25,0)$) node[right, midway] {\small avbrott};
	
	\end{tikzpicture}
	
\end{document}  
}
	\caption{Översikt av systemet \label{system}}	
\end{figure}


\subsection{Beskrivning av systemet}
Roboten navigerar själva banan med hjälp av flertal sensorer, dessa finns specificerade i \mbox{kapitel \ref{sec:sensormodul}}. En regleringsmodell ser till att roboten färdas i mitten av korridorerna och kan ta svängar utan att stöta mot väggar. Under färden ska roboten autonomt kartlägga och finna den kortaste vägen mellan ingången och den nödställde. Ifall roboten även är uppkopplad mot datormodulen ska mjukvara på datorn successivt rita upp en karta och kunna presentera utvalda mätvärden i realtid. När den kortaste vägen är funnen ska roboten använda denna rutt för att förse den nödställde med en förnödenhet. Förnödenheten transporteras med hjälp av robotens gripklo.

\subsection{Delsystem}
Delsystemen och dess ingående komponenter finns beskrivna nedan.
\begin{description}
	\item[Delsystem 1 - Huvudmodul] \hfill \\
	Mikroprocessor \\
	Brytare för autonomt/manuellt-läge \\
	Bluetooth\textsuperscript{\circledR} sändare/mottagare
	\item[Delsystem 2 - Sensormodul] \hfill \\
	Mikroprocessor \\
	Sensorer
	\item[Delsystem 3 - Styrmodul] \hfill \\
	Mikroprocessor \\
	Motorer till hjul och gripklo \\
	LCD-display
	\item[Delsystem 4 - Datormodul] \hfill \\
	Dator \\
	GUI \\
	Bluetooth\textsuperscript{\circledR} sändare/mottagare
\end{description}

Ovanstående mikroprocessorer är av typ \verb+ATmega**+ där versionen bestäms i ett senare skede. Sensorerna består av avståndsmätare, vinkelhastighets-sensorer och identifierare av nödställd. 
\subsection{Kommunikation mellan delsystem}
Kommunikationen mellan mikroprocessorerna sker med hjälp av en I\textsuperscript{2}C-buss enligt figur \ref{communication}. Mellan huvudmodulen och datormodulen sker kommunikationen via Bluetooth\textsuperscript{\circledR}.

\begin{figure}[htbp]
\noindent\resizebox{.97\textwidth}{!}{
	\documentclass[border=20pt]{standalone}
\usepackage{tikz}
\usetikzlibrary{positioning}
\usetikzlibrary{calc}
\usetikzlibrary{decorations.pathmorphing}
\usepackage{amssymb}
\usetikzlibrary{shapes,arrows}

\begin{document}
	\begin{tikzpicture}[scale=1]
		
		\tikzset{every node/.style={thick, draw=black, align=center, minimum height=40pt, text width=100pt, minimum width=100pt}}
		\node(datormodul) {Datormodul};
		\node[right=10pt of datormodul,minimum height=20pt, minimum width=10pt,text width=10pt] (bt1) {\includegraphics{bluetooth}};
		
		\node[right=40pt of bt1,minimum height=20pt, minimum width=10pt,text width=10pt] (bt2) {\includegraphics{bluetooth}};
		
		\node[right=10pt of bt2] 			(huvudmodul) 	{Huvudmodul};
		\node[below=-10pt of huvudmodul,draw=none] (master) {\textit{master}};
		\node[right=10pt of huvudmodul] 		(sensormodul) 	{Sensormodul};
		\node[below=-10pt of sensormodul,draw=none] (slave1) {\textit{slav}};
		\node[right=10pt of sensormodul] 	(styrmodul) 		{Styrmodul};
		\node[below=-10pt of styrmodul,draw=none] (slave2) {\textit{slav}};
		
		\coordinate (sclStart) 	at ($(huvudmodul.north west) + (0,20pt)$);
		\coordinate (sclEnd)		at ($(styrmodul.north east)  + (0,20pt)$);
		
		\coordinate (sdaStart)  at ($(sclStart) + (0,20pt)$);
		\coordinate (sdaEnd)		at ($(sclEnd)	+ (0,20pt)$);
		
		\coordinate (vddStart)  at ($(sdaStart) + (0,40pt)$);
		\coordinate (vddEnd)		at ($(sdaEnd)	+ (0,40pt)$);
		
		\draw[thick] (sclStart) -- (sclEnd) node [right,draw=none,text width=0,minimum width=0] {SCL};
		\draw[thick] (sdaStart) -- (sdaEnd) node [right,draw=none,text width=0,minimum width=0] {SDA};
		\draw[thick] (vddStart) -- (vddEnd) node [right,draw=none,text width=0,minimum width=0] {$V_{dd}$};
		
		\draw[thick] (datormodul.east) -- (bt1.west);
		
		\draw[thick, ->,line join=round,decorate, decoration={
    												snake,
    												segment length=5,
    												amplitude=1,
    												post=lineto,
    												post length=1pt}] 
    		($(bt1.east) + (5pt,5pt)$) -- ($(bt2.west) + (-5pt,5pt)$);
    		
    	\draw[thick, ->,line join=round,decorate, decoration={
    												snake,
    												segment length=5,
    												amplitude=1,
    												post=lineto,
    												post length=1pt}] 
    		 ($(bt2.west) + (-5pt,-5pt)$) -- ($(bt1.east) + (5pt,-5pt)$);
    		 
    	\draw[thick] (bt2.east) -- (huvudmodul.west);
    	
    	\draw[thick,fill=black] ($(huvudmodul.north) + (-10pt,0)$) -- ($(huvudmodul.north) + (-10pt,20pt)$) circle [radius=2pt];
    	\draw[thick,fill=black] ($(huvudmodul.north) + (+10pt,0)$) -- ($(huvudmodul.north) + (+10pt,40pt)$) circle [radius=2pt];
    	
    	\draw[thick,fill=black] ($(sensormodul.north) + (-10pt,0)$) -- ($(sensormodul.north) + (-10pt,20pt)$) circle [radius=2pt];
    	\draw[thick,fill=black] ($(sensormodul.north) + (+10pt,0)$) -- ($(sensormodul.north) + (+10pt,40pt)$) circle [radius=2pt];
    	
    	\draw[thick,fill=black] ($(styrmodul.north) + (-10pt,0)$) -- ($(styrmodul.north) + (-10pt,20pt)$) circle [radius=2pt];
    	\draw[thick,fill=black] ($(styrmodul.north) + (+10pt,0)$) -- ($(styrmodul.north) + (+10pt,40pt)$) circle [radius=2pt];
    	
    	\draw[thick,fill=black] ($(sensormodul.north east) + (-5pt,20pt)$) circle [radius=2pt] -- ($(sensormodul.north east) + (-5pt,80pt)$) circle [radius=2pt];
    	\draw[thick,fill=black] ($(styrmodul.north west) + (5pt,40pt)$) circle [radius=2pt] -- ($(styrmodul.north west) + (5pt,80pt)$) circle [radius=2pt];
    	\draw[thick,draw=black,fill=white] ($(styrmodul.north west) + (2pt,50pt)$) rectangle ($(styrmodul.north west) + (8pt,70pt)$);
		\draw[thick,draw=black,fill=white] ($(sensormodul.north east) + (-8pt,50pt)$) rectangle ($(sensormodul.north east) + (-2pt,70pt)$) node [below left=-10pt and 20pt,draw=none,minimum width=0, text width = 0pt] {$R_p$};
		
		\draw[thick, ->] ($(sensormodul.south) + (+40pt,0)$) -- ($(sensormodul.south) + (40pt,-30pt)$) -- ($(huvudmodul.south) + (40pt,-30pt)$) -- ($(huvudmodul.south) + (40pt,0)$);
		\draw[thick, ->] ($(styrmodul.south) + (+40pt,0)$) -- ($(styrmodul.south) + (40pt,-40pt)$) -- node [midway,below,minimum height=20pt, draw=none] {Avbrott}($(huvudmodul.south) + (30pt,-40pt)$) -- ($(huvudmodul.south) + (30pt,0)$);
	\end{tikzpicture}
\end{document}	}
	\caption{Intermodulär kommunikation \label{communication}}
\end{figure}

\pagebreak
\section{Delsystem 1 - Huvudmodul}
\label{sec:huvudmodul}
text
\subsection{LCD-display}
text
\subsection{Brytare}
% Typ auto/manuell styrning

\pagebreak
\section{Delsystem 2 - Sensormodul}
\label{sec:sensormodul}
text
\subsection{Sensorer}
text

\pagebreak
\section{Delsystem 3 - Styrmodul}
text
\subsection{Motorer}
text
\subsection{Gripklo}
text

\pagebreak
\section{Delsystem 4 - Datormodul}
text
\subsection{Mjukvara}
text
\subsection{Gränssnitt}
text

\setcounter{secnumdepth}{0}
\pagebreak
\begin{thebibliography}{9}

\bibitem{krav}
  Grupp 4,
  \emph{Kravspecifikation 1.0},
  2016-02-03
  
\bibitem{banspec}
	Grupp 1-4,
	\emph{Banspecifikation 1.0}
	
\bibitem{tavling}
	Grupp 1-4,
	\emph{Tävlingsregler 1.0}

\end{thebibliography}

\setcounter{secnumdepth}{2}

\end{flushleft}

\end{document}